\documentclass[pdflatex,sn-mathphys-num]{sn-jnl}% Math and Physical Sciences Numbered Reference Style

\usepackage{graphicx}%
\usepackage{multirow}%
\usepackage{amsmath,amssymb,amsfonts}%
\usepackage{amsthm}%
\usepackage{mathrsfs}%
\usepackage[title]{appendix}%
\usepackage{xcolor}%
\usepackage{textcomp}%
\usepackage{manyfoot}%
\usepackage{booktabs}%
\usepackage{algorithm}%
\usepackage{algorithmicx}%
\usepackage{algpseudocode}%
\usepackage{listings}%

\providecommand{\code}[1]{\texttt{#1}}% used in Section 4
\algrenewcommand\algorithmicrequire{\textbf{Input:}}
\algrenewcommand\algorithmicensure{\textbf{Output:}}

\theoremstyle{thmstyleone}%
\newtheorem{theorem}{Theorem}%
\newtheorem{proposition}[theorem]{Proposition}%

\theoremstyle{thmstyletwo}%
\newtheorem{example}{Example}%
\newtheorem{remark}{Remark}%

\theoremstyle{thmstylethree}%
\newtheorem{definition}{Definition}%

\raggedbottom

\begin{document}

\title[Adaptive Identification and Tracking]{Adaptive Candidate Routing for Distributed Occupant Recognition Across Multiple Building IoT-enabled Smart Spaces.}

\author*[1,2]{\fnm{First} \sur{Author}}\email{iauthor@gmail.com}

\author[2,3]{\fnm{Second} \sur{Author}}\email{iiauthor@gmail.com}
\equalcont{These authors contributed equally to this work.}

\author[1,2]{\fnm{Third} \sur{Author}}\email{iiiauthor@gmail.com}
\equalcont{These authors contributed equally to this work.}

\affil*[1]{\orgdiv{Department}, \orgname{Organization}, \orgaddress{\street{Street}, \city{City}, \postcode{100190}, \state{State}, \country{Country}}}

\affil[2]{\orgdiv{Department}, \orgname{Organization}, \orgaddress{\street{Street}, \city{City}, \postcode{10587}, \state{State}, \country{Country}}}

\affil[3]{\orgdiv{Department}, \orgname{Organization}, \orgaddress{\street{Street}, \city{City}, \postcode{610101}, \state{State}, \country{Country}}}

\abstract{In a distributed state transition system (DSTS) for multi-building camera networks, a face embedding captured in one building must be traced to the building-specific state transition system (BSTS) in which its occupant is registered, before the recognition can be translated into an event and the states updated. A centralized master gallery makes this easy but keeps sensitive biometric data in one place. We propose a routing strategy that avoids it: each BSTS publishes one Bloom filter built from locality-sensitive hashing (LSH) codes of the face embeddings of its registered occupants, a captured face embedding is checked against all filters, and only the best-scoring BSTSs verify it against their own reference face embeddings. On 500 CASIA-WebFace identities in 10 buildings, averaged over 20 random building assignments, the default setting identifies 92.37\% of test face embeddings correctly, against 95.47\% for a centralized search with the same vote rule; an adaptive rule that asks extra BSTSs only when the ranking is unclear reaches 92.44\% and contacts 2.54 BSTSs on average. LSH code collisions raise the real false-positive rate of the filters to 3.54\% instead of the configured 1\%. Over 2, 5 and 10 BSTSs, the adaptive rule stays within 0.5 points of the centralized search at 2 and 5 BSTSs and 2.9 points below it at 10, and someone who holds the filters can tell which building an occupant is registered in (AUC 0.93). \textcolor{blue}{Costs are computed, not measured; encryption and unknown individuals are not evaluated}.}

\keywords{Occupant identification and tracking, Distributed state transition system, Smart Spaces, Locality-sensitive hashing, Bloom filter, Privacy in data analytics}

\maketitle

\section{Introduction}\label{sec1}

Large camera networks that monitor several buildings can be abstracted as a distributed state transition system (DSTS), composed of independent building-specific state transition systems (BSTSs) in which recognition events drive state transitions~\cite{mohan2020}. Every occupant is registered in exactly one BSTS, and an occupant seen in another building is a visitor there. When a face embedding is captured in one building, the system must therefore find the BSTS in which its occupant is registered before the recognition can be translated into an event and the states updated. The DSTS framework resolves such visitors against a centralized master list of registered occupants~\cite{mohan2020}, that is, one central gallery with the face embeddings of everyone. This is accurate, but it stores sensitive biometric data in a single place, an attractive target that may conflict with data-minimization rules. The other extreme, a request to every BSTS for every captured face embedding, exposes every captured face embedding to every building.

We study the step between these two: a routing step that decides which few BSTSs should be asked. Each BSTS publishes one Bloom filter~\cite{bloom1970} built from LSH codes of the face embeddings of its registered occupants. A captured face embedding is checked against all filters, and only the best-scoring BSTSs receive it. Reference face embeddings are not published. Face matching is approximate, so the summary cannot hold exact keys, as in web cache sharing~\cite{fan2000}; it holds LSH codes instead.

\emph{Scope.} This paper is about the routing step: its accuracy, the number of BSTSs it contacts, the false positives of the filters, and its behaviour with 2, 5 and 10 BSTSs. All evaluated face embeddings belong to registered occupants; unknown individuals, who are registered in no BSTS, are left for future work. We do not implement any cryptographic protection (Section~\ref{sec:model}), and costs are computed from message sizes, not measured in a distributed run.

Our contributions are:
\begin{itemize}
\item a routing design with one Bloom filter of LSH codes per BSTS, and a way to choose the code length $K$ and the number of codes $L$ from measured angle distributions;
\item measurements on 500 identities and 20 building assignments: accuracy against a centralized search, the gain from multi-probe hashing, and a real false-positive rate of 3.54\% that we trace to collisions between LSH codes;
\item a rule that asks extra BSTSs only when the two best scores are close, which is as accurate as asking all shortlisted BSTSs and contacts about half as many;
\item an evaluation with 2, 5 and 10 BSTSs, an analytic cost comparison, and a measurement of what the published filters reveal about the building in which an occupant is registered.
\end{itemize}

Section~\ref{sec:related} reviews studies on multi-building occupant tracking, covering centralized and distributed approaches, and the work on Bloom filters and routing in distributed search that our design builds on. Section~\ref{sec:method} develops the routing method, from the system and threat model through LSH codes and per-BSTS Bloom filters to adaptive routing that ends in event and state generation. Section~\ref{sec:setup} describes the data, parameters and figures of merit; Section~\ref{sec:results} reports the measurements; and Section~\ref{sec:limits} states what the study does not cover and concludes.

%=====================================================================
\section{Related Work}
\label{sec:related}
Existing studies in the domain of indoor occupant tracking focus on camera-based~\cite{mohan2020} or RF-signal-based approaches~\cite{WZSYGWLLHSC2025}. These works highlight computational efficiency and privacy preservation as two major challenges in occupant identification and tracking across multi-building indoor smart spaces, especially when camera networks are employed~\cite{KKZWKMW2025}. Consequently, we review key studies that address privacy concerns and computational requirements associated with occupant tracking in IoT-enabled multi-building indoor smart spaces.

Occupant identification and tracking are essential to context-aware applications in smart-space environments, which become more intelligent and context-aware through the integration of IoT-enabled sensing devices~\cite{KZGS2024}. Facial biometrics are particularly well suited for camera-based human identification and tracking in smart spaces because they provide distinctive identity information and are easily captured remotely. Unlike many other biometric modalities, face recognition can operate unobtrusively without requiring the individual to stop, actively cooperate, or explicitly present an identifying credential.

\emph{State transition models.} Occupant tracking has been modelled with state transition systems whose events abstract the recognition step~\cite{MJG2011}. The DSTS extends this to multi-building camera networks as independent BSTSs and resolves a visitor against a centralized master list of registered occupants~\cite{mohan2020}; we replace that lookup by routing. Further work covers video synopsis for camera-based surveillance analysis~\cite{NGA2022} and quantitative metrics for smart spaces~\cite{MMJ2025}.

\emph{Template protection and Bloom filters.} Rathgeb et al.~\cite{rathgeb2014} apply Bloom filters to iris templates. Closest to our work, Marzani et al.~\cite{marzani2026} turn face features into identifiers with a fuzzy extractor and store them in homomorphically encrypted Bloom filters to count people across locations; their goal is counting by exact identifier membership, ours is routing a face embedding by approximate similarity (multi-probe LSH codes) to the BSTS of its registered occupant.

\emph{Routing in distributed search.} Summary Cache~\cite{fan2000} lets web proxies exchange Bloom filters of their contents and contact only matching proxies; we follow this idea with LSH codes in place of exact keys. For vector data, Gottesb{\"u}ren et al.~\cite{gottesbueren2024} study clustering- and LSH-based methods that send a request only to shards likely to hold its neighbours, and Semord~\cite{wang2026semord} places similar vectors close together in a decentralized overlay with a learned hash; its peers exchange summaries that include representative vectors, and it is evaluated on generic benchmarks, not on face embeddings. These systems route by vector placement, centroids or learned structure, which requires moving data or sharing vector summaries; we let each BSTS publish a small Bloom filter and keep its reference face embeddings. To our knowledge, per-BSTS Bloom filters of LSH codes have not been measured for this task before.

%=====================================================================
\section{Methodology}
\label{sec:method}

The method has two phases: each BSTS builds its Bloom filter once (Algorithm~\ref{alg:build}), and every captured face embedding is then routed, verified, and translated into a recognition event (Algorithm~\ref{alg:adaptive-routing}). We first state the setting and what is protected, then define the quantities that the algorithms use.

\subsection{System and Threat Model}
\label{sec:model}

We build on the distributed state transition system (DSTS) of~\cite{mohan2020}, which abstracts a multi-building camera network as independent building-specific state transition systems (BSTSs). Each BSTS$_b$ has a set of states, a set of events and a transition function $\Delta^b:S^b\times E^b\to S^b$; a state records the probabilities of occupants being present in the zones of the building~\cite{mohan2020,MJG2011}. BSTS$_b$ stores the reference face embeddings of its $N$ registered occupants; an occupant registered in another BSTS is a visitor there. A face captured at event instance $k$ in a zone of building $q$ is converted into a face embedding. Recognition of this face embedding is translated into a recognition event $e_k$, and the transition function derives the next state $S_k$ from the current state $S_{k-1}$. When the occupant is registered in another BSTS, the DSTS looks the face embedding up in a centralized master list of all registered occupants~\cite{mohan2020}; this paper replaces that lookup by the routing step below and evaluates the identification step only. Every occupant is registered in exactly one BSTS. Each BSTS builds a Bloom filter from the Locality-Sensitive Hashing (LSH) codes of the face embeddings of its registered occupants; the filters shortlist candidate BSTSs, which verify the captured face embedding against their stored reference face embeddings. Reference face embeddings are never exchanged, although the captured face embedding is sent to shortlisted BSTSs. BSTSs are assumed to execute the operations correctly but may try to infer sensitive information from what they receive; published filters, public hash parameters and captured face embeddings may reveal building membership (Section~\ref{sec:leak}). Authenticated, encrypted communication is assumed but not implemented. Face guessing, external information, repeated captures, colluding or malicious BSTSs, and traffic analysis are out of scope.

\subsection{Face Embedding and LSH Codes}
For the face $o_k$ captured at event instance $k$, SCRFD-10GF~\cite{guo2022} detects the face and pretrained ArcFace (InsightFace \code{buffalo\_l})~\cite{deng2019} produces a 512-dimensional raw face embedding $\tilde{\mathbf{x}}_k$; reference face embeddings are produced in the same way. We first unit-normalize and then subtract the reference-only population mean $\bar{\mathbf{x}}$ and renormalize, giving the face embedding $\mathbf{p}$ used for routing:
\begin{equation}
\mathbf{x}_k=\frac{\tilde{\mathbf{x}}_k}{\|\tilde{\mathbf{x}}_k\|_2},\qquad
\mathbf{p}=\frac{\mathbf{x}_k-\bar{\mathbf{x}}}{\|\mathbf{x}_k-\bar{\mathbf{x}}\|_2}.
\label{eq:preprocess}
\end{equation}
The same mean and transformation are used for enrollment and identification. Let $K$ be the number of bits per code and $L$ the number of slots. Using $KL$ fixed-seed Gaussian hyperplanes $\mathbf{h}_{\ell,i}$, define $b_{\ell,i}=\mathbb{1}[\mathbf{p}\cdot\mathbf{h}_{\ell,i}\geq0]$ and the slot-tagged item
\begin{equation}
c_\ell=\sum_{i=0}^{K-1}b_{\ell,i}2^i,\qquad s_\ell=\ell2^K+c_\ell.
\label{eq:lsh-code}
\end{equation}
For angle $\theta$ between two unit face embeddings, a bit matches with probability $p_{\rm bit}=1-\theta/180$~\cite{charikar2002}.

\subsection{Parameter Selection}
Averaging $p_{\rm bit}^K$ over reference-only genuine and impostor pairs of face embeddings gives $p_{\rm code}^{(G)}$ and $p_{\rm code}^{(I)}$. Assuming independent slots,
\begin{equation}
\mathrm{Recall}=1-(1-p_{\rm code}^{(G)})^L,\quad
\mathrm{FPR}_{\rm imp}=1-(1-p_{\rm code}^{(I)})^L.
\label{eq:lsh-recall}
\end{equation}
For a recall target $\rho$, set $L=\left\lceil\ln(1-\rho)/\ln(1-p_{\rm code}^{(G)})\right\rceil$. For each candidate $K$ we compute $L$, reject values with predicted $\mathrm{FPR}_{\rm imp}>\varphi$ for a chosen bound $\varphi$, and keep the pair with the smallest $KL$.

\subsection{Per-BSTS Bloom Filters}
Each BSTS publishes a Bloom filter but retains its reference face embeddings. For registered occupant $u$ in BSTS$_b$, the normalized mean of the $R$ reference face embeddings forms centroid $\mathbf{z}_{b,u}=\operatorname{Normalize}(R^{-1}\sum_{r=1}^{R}\mathbf{e}_{b,u,r})$, which is encoded into $L$ items using Equation~\ref{eq:lsh-code}. With $n=NL$ inserted items and target false-positive probability $f^\star$, the filter size and number of hash positions are
\begin{equation}
m=\left\lceil\frac{-n\ln f^\star}{(\ln2)^2}\right\rceil,\qquad
\kappa=\operatorname{round}\!\left(\frac{m}{n}\ln2\right).
\label{eq:bloom-size}
\end{equation}
We use Kirsch--Mitzenmacher double hashing~\cite{kirsch2006}: $h_1=\mathrm{SM64}(s\oplus(2\sigma+1))$, $h_2=\mathrm{SM64}(s\oplus(2\sigma+2))\vee1$, and $p_j=((h_1+jh_2)\bmod2^{64})\bmod m$ for $j=0,\ldots,\kappa-1$, where $\sigma$ is a shared salt and SM64 is the SplitMix64 finalizer~\cite{steele2014}. A check is positive when all $\kappa$ positions are set.

\begin{algorithm}[!t]
\footnotesize
\caption{Building the published Bloom filter at BSTS$_b$}
\label{alg:build}
\begin{algorithmic}[1]
\Require Raw reference face embeddings $\tilde{\mathbf{x}}_{b,u,r}$ of the $N$ registered occupants of BSTS$_b$ ($R$ per occupant);
\Statex population mean $\bar{\mathbf{x}}$; shared hyperplanes $\mathcal{H}$ and salt $\sigma$; parameters $K,L,f^\star$.
\Ensure Published Bloom filter $F_b$; private reference face embeddings $\mathcal{E}_b$.
\State $n\gets NL$; derive $m$ and $\kappa$ \Comment{Equation~\ref{eq:bloom-size}}
\State $F_b\gets$ all-zero bit array of length $m$
\For{$u=1$ to $N$}
  \For{$r=1$ to $R$}
    \State $\mathbf{e}_{b,u,r}\gets$ preprocessed $\tilde{\mathbf{x}}_{b,u,r}$ \Comment{Equation~\ref{eq:preprocess}}
  \EndFor
  \State $\mathbf{z}_{b,u}\gets\operatorname{Normalize}\bigl(R^{-1}\sum_{r}\mathbf{e}_{b,u,r}\bigr)$ \Comment{occupant centroid}
  \For{$\ell=0$ to $L-1$}
    \State $s_\ell\gets$ slot-tagged item of $\mathbf{z}_{b,u}$ \Comment{Equation~\ref{eq:lsh-code}}
    \State set the $\kappa$ bits $p_0,\dots,p_{\kappa-1}$ of item $s_\ell$ in $F_b$ \Comment{double hashing}
  \EndFor
\EndFor
\State $\mathcal{E}_b\gets\{\mathbf{e}_{b,u,r}\}$, kept private; publish $F_b$
\end{algorithmic}
\end{algorithm}

\subsection{Routing, Verification, and Event Generation}\label{sec:escalation}
For each slot, multi-probe routing~\cite{lv2007} checks the original code and $P$ single-bit variants obtained by flipping the bits with the smallest margins $|\mathbf{p}\cdot\mathbf{h}_{\ell,i}|$:
\begin{equation}
s'_{\ell,j}=\ell2^K+\bigl(c_\ell\mathbin{\mathrm{XOR}}2^{w_{\ell,j}}\bigr),\quad j=1,\ldots,P,
\label{eq:multi-probe}
\end{equation}
where $w_{\ell,j}$ is the index of the $j$-th smallest-margin bit. BSTS$_b$ receives score $A_b$, the number of slots with at least one positive check; the shortlist $\mathcal{C}$ contains the top $\min(S,B)$ BSTSs by score (ties by index). In a shortlisted BSTS, each reference face embedding votes for its registered occupant when its angle to the captured face embedding is at most the accept angle $\tau$, the cutoff that minimizes the sum of the genuine false-reject rate and the impostor false-accept rate on reference face embeddings; the local winner is the registered occupant with most votes (ties by smallest reference angle), and acceptance requires at least $V$ of $R$ votes. The adaptive policy verifies all shortlisted BSTSs when $g=A_{(1)}-A_{(2)}\leq g_0$ and selects by vote count, minimum reference angle, then routing rank; otherwise it stops at the first local acceptance in routing order. An accepted occupant is translated into the recognition event $e_k$ of the BSTS in which it is registered, and the transition function updates the state. If none is accepted, the system abstains, no event is generated and the state is left unchanged. The default all-shortlist rule used in the multi-probe comparison instead compares votes across all registered occupants of the shortlisted BSTSs, breaking ties by minimum reference angle.

\begin{algorithm}[!t]
\footnotesize
\caption{Adaptive candidate routing, event generation and state generation}
\label{alg:adaptive-routing}
\begin{algorithmic}[1]
\Require Face embedding $\tilde{\mathbf{x}}_k$ captured at event instance $k$ in zone $j$ of building $q$;
\Statex published filters $\{F_b\}_{b=1}^{B}$; private reference face embeddings $\{\mathcal{E}_b\}_{b=1}^{B}$ held by the BSTSs;
\Statex mean $\bar{\mathbf{x}}$, hyperplanes $\mathcal{H}$, parameters $K,L,P,S,V,\tau,g_0$;
\Statex current state $S_{k-1}$.
\Ensure Recognised occupant $a^\star$ registered in BSTS$_{b^\star}$ (or $\mathrm{ABSTAIN}$), event $e_k$ and next state $S_k$.
\State $\mathbf{p}\gets$ preprocessed $\tilde{\mathbf{x}}_k$ \Comment{Equation~\ref{eq:preprocess}}
\For{$\ell=0$ to $L-1$}
  \State compute $s_\ell$ and the $P$ probe items $s'_{\ell,1},\dots,s'_{\ell,P}$ \Comment{Eqs.~\ref{eq:lsh-code}, \ref{eq:multi-probe}}
\EndFor
\For{$b=1$ to $B$}
  \State check the original and probe items of every slot against $F_b$
  \State $A_b\gets$ number of slots with at least one positive check \Comment{each slot once}
\EndFor
\State $\mathcal{C}\gets$ top $\min(S,B)$ BSTSs by decreasing $A_b$ \Comment{ties by index}
\State $g\gets A_{(1)}-A_{(2)}$ if $B\geq2$, else $g\gets+\infty$
\State $a^\star\gets\mathrm{ABSTAIN}$
\If{$g\leq g_0$} \Comment{ranking unclear: verify all shortlisted BSTSs}
  \ForAll{$b\in\mathcal{C}$}
    \State $(u_b,v_b)\gets\textsc{LocalVerify}(b,\mathbf{p})$
  \EndFor
  \State $\mathcal{A}\gets\{b\in\mathcal{C}:v_b\geq V\}$
  \If{$\mathcal{A}\neq\emptyset$} \State $(a^\star,b^\star)\gets(u_b,b)$ for the $b\in\mathcal{A}$ with most votes \Comment{ties: angle, then rank} \EndIf
\Else \Comment{ranking clear: sequential first-match}
  \ForAll{$b\in\mathcal{C}$ in routing order}
    \State $(u_b,v_b)\gets\textsc{LocalVerify}(b,\mathbf{p})$
    \If{$v_b\geq V$} \State $(a^\star,b^\star)\gets(u_b,b)$; \textbf{break} \EndIf
  \EndFor
\EndIf
\State $e_k\gets\textit{event\_gen}(a^\star,b^\star,k,q,j)$ \Comment{recognition translated into an event; none on $\mathrm{ABSTAIN}$}
\State $S_k\gets\textit{state\_gen}(S_{k-1},e_k)$ \Comment{transition function $\Delta$; $S_k=S_{k-1}$ on $\mathrm{ABSTAIN}$}
\State \Return $a^\star,e_k,S_k$
\Function{LocalVerify}{$b,\mathbf{p}$} \Comment{runs in BSTS$_b$ on its private reference face embeddings}
  \ForAll{registered occupant $u$ of BSTS$_b$}
    \State $v_u\gets|\{r:\angle(\mathbf{p},\mathbf{e}_{b,u,r})\leq\tau\}|$
  \EndFor
  \State \Return $u^\star=\arg\max_u v_u$ (ties by smallest reference angle) and $v_{u^\star}$
\EndFunction
\end{algorithmic}
\end{algorithm}

\section{Experimental Setup}
\label{sec:setup}

We use the face embeddings of 500 identities of CASIA-WebFace~\cite{yi2014}, registered in $B=10$ BSTSs of $N=50$ occupants each. Each occupant has $R=20$ reference face embeddings and 20 held-out test face embeddings, giving 20,000 face embeddings. Every one of the 10,000 test face embeddings is processed as one event instance: it is routed over all $B$ BSTSs and its recognition is translated into an event. The evaluation is \emph{closed-set}: every test face embedding belongs to an occupant registered in exactly one BSTS. We use $K=11$, $L=122$, $P=3$, $S=5$, $V=12$ and $f^{\star}=0.01$. With these settings each filter stores $n=NL=6{,}100$ items and has $m=58{,}469$ bits and $\kappa=7$ hash positions (about 7.3\,KB). For parameter selection (Section~\ref{sec:method}) we use a recall target $\rho=0.90$, an impostor-FPR bound $\varphi=0.10$ and $K\in\{4,\ldots,15\}$; the model selects $(K,L)=(10,88)$, whereas the main experiments use $(11,122)$ because it performed better on test accuracy (Section~\ref{sec:k1011}).

\emph{Use of test data.} The mean face embedding, the accept angle and the angle statistics behind $K$ and $L$ are estimated from reference face embeddings only (genuine pairs are two reference face embeddings of the same occupant, impostor pairs those of different occupants); this gives $\tau=80.43^{\circ}$. Using test face embeddings for the mean and $\tau$ changed correct-identification rates by at most 0.06 points. Since $K=11$ and $P=3$ were chosen after looking at test accuracy, we also report $K=10$, the model's choice, and the full sweep over $P$ (Section~\ref{sec:k1011}).

\emph{Seeds.} Each result is repeated for 20 seeds. A seed changes the assignment of the 500 identities to BSTSs, the LSH hyperplanes and the Bloom filter salt; the face embeddings, their mean and the accept angle stay fixed, so the centralized baseline is a single number. We report the mean over seeds with a 95\% confidence half-width ($\pm$, Student-$t$) and, for comparisons, the paired difference over seeds.

\emph{Metrics.} Each test face embedding results in exactly one of four recognition outcomes: \emph{correct} (the event names the right registered occupant); \emph{misidentified} (a wrong occupant is accepted, the most serious error); \emph{abstained-reject} (the true BSTS was shortlisted but nobody reached $V$ votes); and \emph{abstained-routing} (the true BSTS was not shortlisted and nobody reached $V$ votes). The centralized baseline gives 95.47\% correct, 1.41\% misidentified and 3.12\% abstained.

%=====================================================================
\section{Results}
\label{sec:results}

\subsection{Choice of $K$ and $L$}
\label{sec:k1011}

The analytic model, which uses only the angle arrays, picks $K=10$, $L=88$, but its mean-angle step underestimates the genuine code probability ($0.0115$ against $0.0187$ at $K=11$). In the full pipeline at $P=3$, $K=11$, $L=122$ is more accurate by $0.52 \pm 0.13$ points (paired, 20 of 20 seeds): 92.37\% against 91.85\%, mainly because routing failures fall from 4.88\% to 4.46\%. We use $K=11$ and note that this choice was made after seeing test accuracy.

\subsection{Multi-Probe and Decentralized versus Centralized Identification}

Table~\ref{tab:probe} varies $P$ at $B=10$ with the default rule. Without probing, $8.07 \pm 0.11\%$ of face embeddings are lost to routing misses; $P=3$ reduces this to $4.46 \pm 0.08\%$, raises correct identification from $88.28\%$ to $92.37\%$ (paired gain $4.09 \pm 0.15$ points, 20 of 20 seeds) and lowers misidentification from 1.91\% to 1.42\%. The gains shrink as $P$ grows while lookups grow linearly, so $P=3$ is a compromise on cost. The decentralized system is then $3.10 \pm 0.08$ points behind the centralized baseline, and 4.46\% of face embeddings are routing failures. Before multi-probe the gap was 7.19 points, so multi-probe closes 57\% of it. Misidentification (1.42\% against 1.41\%) hardly changes, so the loss comes from missed matches, not from more wrong accepts.

\begin{table}[!t]
\caption{Multi-probe sweep at $B=10$ ($K=11$, $L=122$, default rule).
Results are reported as the mean over 20 seeds ($\pm$ 95\% CI), in \%.
The number of items looked up per filter and face embedding is $L(1+P)$.}
\label{tab:probe}

\begin{tabular*}{\textwidth}
{@{\extracolsep\fill}cccccc}
\toprule
$P$
& Correct & Misid. & Abst.-reject & Abst.-routing & Items \\
\midrule

0 & $88.28 \pm 0.13$ & $1.91 \pm 0.04$ & $1.74 \pm 0.03$ & $8.07 \pm 0.11$ & 122 \\

3 & $92.37 \pm 0.08$ & $1.42 \pm 0.03$ & $1.75 \pm 0.04$ & $4.46 \pm 0.08$ & 488 \\

7 & $92.89 \pm 0.07$ & $1.38 \pm 0.03$ & $1.76 \pm 0.05$ & $3.97 \pm 0.08$ & 976 \\
\botrule
\end{tabular*}

\end{table}

\subsection{Effective False-Positive Rate}
\label{sec:fpr}

The 1\% target describes the bit array for items that are independent of the stored ones. Testing the enrolled items of the \emph{other nine} BSTSs against each filter gives a real per-item rate of $3.54 \pm 0.02\%$ over the 20 seeds. Comparing the integer codes directly, $2.61 \pm 0.02\%$ of those items already equal an enrolled item in the same slot, close to the 2.41\% expected if the codes of 50 occupants were evenly spread over $2^{11}$ values, whereas the Bloom array alone, tested on random integers, gives $0.94 \pm 0.01\%$ at 51.3\% fill. The extra false positives therefore come from the LSH codes, not from the Bloom array. A larger $K$ lowers the rate but $L$ must grow: at $K=13$, $L=230$ (13.8\,KB per filter, 1,610 lookups per filter and face embedding against 488) the real rate falls to 1.67\% and, with $P=6$, accuracy is 93.10\%, $0.73 \pm 0.11$ points above the default.

\subsection{Verification Policies}
\label{sec:adaptive}

Figure~\ref{fig:results}(a) shows the policies of Section~\ref{sec:escalation} and the fixed shortlist sizes; the figures in this paragraph are for the 10,000 test face embeddings of the main experiment. Sequential first-match raises misidentification to 2.17\%, because a wrong BSTS ranked above the true one sometimes confirms first. Asking all five BSTSs lowers it to 1.45\% and gains $0.71 \pm 0.03$ points but contacts 5 BSTSs per face embedding. The adaptive rule with $g_0=5$ is slightly more accurate than asking all five (92.44\% against 92.34\%, $+0.10 \pm 0.02$ points, paired, 20 of 20 seeds) while contacting 2.54 BSTSs on average; 37.6\% of face embeddings ask all five. All three policies abstain on the same face embeddings, because acceptance depends only on whether some shortlisted BSTS reaches $V$ votes. (Verify-all with this rule is 0.03 points below the default rule of Table~\ref{tab:probe} because of ties between BSTSs.)

\subsection{Cost}
\label{sec:cost}

Costs are computed from message sizes and the measured number of BSTSs asked, not measured in a distributed run. A face embedding has 512 single-precision values (2,048 bytes). Routing replaces a 20.5\,MB master gallery with ten published filters of 7.3\,KB each. Handling one captured face embedding costs 5.2\,KB with the adaptive rule (2.54 BSTSs), against 20.5\,KB for asking all ten BSTSs and 2.0\,KB for a master gallery, plus $B\cdot L(1+P)=4{,}880$ filter lookups, which grow linearly with $B$.

\subsection{What the Filters Reveal}
\label{sec:leak}

Someone who holds all published filters and the public hash parameters can score a face embedding against every filter, as routing does. With $P=3$ this score separates the registered occupants of a BSTS from everyone else with an AUC of 0.93 (0.88 without probing), averaged over BSTSs and 10 seeds, and the top-scoring BSTS is the occupant's own for 79.7\% of face embeddings (66.7\% without probing). The filters therefore leak which building an occupant is registered in; this does not recover a face, but the leak grows with the accuracy of the routing. We did not test attacks with many face embeddings, candidate guessing or outside information, so we make no claim of biometric privacy.

\subsection{Fewer and More BSTSs}
\label{sec:scale}

The corpus has 500 identities, so we split the same 500 registered occupants over $B=2$, 5 and 10 BSTSs of 250, 100 and 50 occupants. Each filter is sized for its own number of occupants. We use 2,000 test face embeddings per seed (4 per occupant) and 10 seeds; this subset gives 92.6\% for the adaptive rule at $B=10$, against 92.44\% on all 10,000 face embeddings. The shortlist holds at most $S=\min(5,B)$ BSTSs, so at $B=5$ it covers every BSTS (recall 100\%), and Figure~\ref{fig:results}(b) therefore reports recall for a shortlist of two. That recall falls from 100\% at $B=2$ to 89.6\% at $B=5$ and 87.0\% at $B=10$, because the shortlist is a smaller share of the BSTSs. Correct identification with the adaptive rule is 95.0\% and 95.3\% at $B=2$ and 5, close to the 95.5\% of the centralized search, and 92.6\% at $B=10$, where the adaptive rule contacts 2.5 BSTSs.

\begin{figure}[!t]
\centering
\includegraphics[width=0.95\textwidth]{results_plots}
\caption{(a) Correct identification against the number of BSTSs contacted at $B=10$: fixed shortlist of $s$ BSTSs, sequential first-match, the adaptive rule and the centralized search. (b) Recall of a shortlist of two and correct identification (adaptive rule and centralized search) for 2, 5 and 10 BSTSs. Mean over 10 seeds with 95\% intervals; 2,000 test face embeddings per seed.}
\label{fig:results}
\end{figure}

%=====================================================================
\section{Limitations and Conclusion}
\label{sec:limits}

\emph{Limitations.} All results use the same 500 identities from one public dataset; the intervals reflect building assignments and hash randomness, not a different population. $K=11$ and $P=3$ were chosen after seeing test accuracy, and alternatives are reported. CASIA-WebFace data come from the web, not from deployed camera networks, so real accuracy is probably lower, and we did not check the overlap between the ArcFace training data and CASIA-WebFace. The largest network has 10 BSTSs. The evaluation is closed-set; a first test with held-out identities showed that the vote rule ($\tau$, $V=12$) can accept a noticeable share of them, so it must be tightened before unknown individuals are handled. We evaluated the identification step only, not the generated events, the mapping of recognition scores to occupant-presence probabilities or the state updates over event sequences. Costs are computed, not measured; encryption, template storage during visits and attacks beyond the membership test of Section~\ref{sec:leak} were not implemented or evaluated. A Bloom filter cannot delete items, so a change in registration needs a rebuild.

\emph{Conclusion.} Per-BSTS Bloom filters of LSH codes can route captured face embeddings to the BSTS of the registered occupant without a centralized master gallery. With multi-probe hashing the default setting identifies 92.37\% of test face embeddings correctly at 10 BSTSs against 95.47\% for a centralized search, and the adaptive rule reaches 92.44\% while contacting about 2.54 BSTSs. The cost is lost recall that grows with the number of BSTSs, and a false-positive rate that LSH collisions raise well above the Bloom filter's design value. The design avoids publishing reference face embeddings and a master gallery, but it leaks which building an occupant is registered in (AUC 0.93). Next steps are a larger and more varied population, unknown individuals with a stricter acceptance rule, and measurements of event generation and state updates over event sequences in a distributed deployment.

\backmatter

\bmhead{Acknowledgement}

The authors gratefully acknowledge Sri Mata Amritanandamayi Devi (Amma), Chancellor, Amrita Vishwa Vidyapeetham, for her inspiration and for providing financial support for the Article Processing Charges (APC) of this publication.

\bibliography{akd}% akd.bib

\end{document}
